\documentclass{article}
\usepackage[utf8]{inputenc}
\usepackage[spanish]{babel}
\usepackage{enumitem}
\usepackage{geometry}
\geometry{a4paper, margin=2.5cm}

\title{ITIL: Information Technology Infrastructure Library}
\author{}
\date{}

\begin{document}
\maketitle

\section{¿Qué es ITIL y cuál es su propósito?}

ITIL (Information Technology Infrastructure Library) es un conjunto de mejores prácticas para la gestión de servicios de TI (ITSM). Fue desarrollado en la década de 1980 por la Central Computer and Telecommunications Agency (CCTA) del gobierno del Reino Unido.

Su propósito es alinear los servicios de TI con las necesidades del negocio, estandarizando procesos y mejorando la calidad, eficiencia y seguridad en la entrega de servicios. ITIL permite a las organizaciones reducir riesgos, optimizar costos, mejorar la satisfacción del cliente y establecer una mejora continua.

\section{Evolución de ITIL y sus versiones principales}

\begin{enumerate}
    \item \textbf{ITIL v1 (1989--1999):} Versión original compuesta por más de 30 libros. Era funcional y carecía de un modelo de proceso central.
    \item \textbf{ITIL v2 (2001):} Consolidó los 30 libros en 7 títulos principales, con un enfoque basado en procesos. Se centró en Service Support y Service Delivery. Se convirtió en el estándar global de ITSM.
    \item \textbf{ITIL v3 (2007):} Introdujo el enfoque de ciclo de vida del servicio con 5 volúmenes: Service Strategy, Service Design, Service Transition, Service Operation y Continual Service Improvement. Incluyó 26 procesos.
    \item \textbf{ITIL 2011:} Actualización de v3 para corregir errores e inconsistencias, mejorando la claridad y alineación entre los volúmenes.
    \item \textbf{ITIL 4 (2019):} Versión actual. Abandona el ciclo de vida y adopta un enfoque holístico basado en el Service Value System (SVS) y las 4 dimensiones. Integra principios de Agile, DevOps y Lean. Se reduce el énfasis en TI hacia una gestión de servicios general.
    \item \textbf{ITIL 5 (2026):} Anunciada en enero de 2026 por PeopleCert como la próxima evolución del marco.
\end{enumerate}

\section{Principios fundamentales de ITIL 4}

ITIL 4 define 7 principios guía que son recomendaciones universales aplicables en cualquier organización:

\begin{enumerate}[label=\textbf{\arabic*.}]
    \item \textbf{Focus on value} --- Todo lo que la organización hace debe generar valor directa o indirectamente para los stakeholders.
    \item \textbf{Start where you are} --- No empezar de cero; evaluar el estado actual y reutilizar lo que ya funciona.
    \item \textbf{Progress iteratively with feedback} --- Avanzar de forma iterativa utilizando retroalimentación para mejorar continuamente.
    \item \textbf{Collaborate and promote visibility} --- Colaborar entre equipos y hacer visibles los procesos y resultados.
    \item \textbf{Think and work holistically} --- Pensar y trabajar de forma integral, considerando el sistema completo.
    \item \textbf{Keep it simple and practical} --- Mantener soluciones simples, prácticas y sin burocracia innecesaria.
    \item \textbf{Optimize and automate} --- Optimizar procesos y automatizar tareas repetitivas cuando sea posible.
\end{enumerate}

\end{document}
